\section{Parametric baseline: additional calibration and diagnostics}\label{sec:supp_parametric}

The main text reports a single ATM term-structure fit to establish the parametric $\psi$ as a baseline. This supplement reports the two earlier-stage parametric calibrations that preceded the neural surrogate: a single-DTE fit on four semiconductor tickers that demonstrated cross-ticker transfer of the smile shape, and the full multi-expiration NVDA calibration from which the main-text figure was drawn, including the parameter values, residual diagnostics, and IV surface reconstruction.

\subsection{Single-DTE four-ticker semiconductor calibration}

The empirical starting point for the parametric form was a set of option chains for four semiconductor stocks: NVDA, AMD, MU, and INTC. All chains were single-date snapshots sourced from barchart.com (October to November 2025) with DTEs ranging from 41 to 80 days, filtered to moneyness $K/S \in [0.70, 1.30]$ and IV $< 200\%$ to exclude illiquid deep ITM/OTM contracts with unreliable IV (Table~\ref{tab:supp_data}). Because each chain was a single-date snapshot, all observations shared one HMM state $s_0$ and market mood $M_0$, so the model IV at a given strike reduced to $\sigma_{\text{model}}(K) = \sqrt{\theta_{\text{base}} \cdot \psi(\text{DTE},\; K/S)}$, where $\theta_{\text{base}} = \theta_{s_0} \cdot (1 + \gamma M_0)$ absorbed the state and mood into a single scalar. We fixed $\beta_1 = 0.20$ (term-structure slope) because single-DTE snapshots could not identify it; the remaining parameters $(\theta_{\text{base}}, \beta_2, \beta_3, \beta_4)$ were fitted by minimizing the mean squared IV error on NVDA using Nelder-Mead optimization. We chose NVDA as the calibration ticker because it had the richest chain (142 contracts after filtering).

\begin{table}[ht]
\centering
\caption{Option chain summary for the single-DTE semiconductor dataset. Contract counts are reported after filtering to the tradeable moneyness range.}\label{tab:supp_data}
\begin{tabular}{lrrrr}
\toprule
Ticker & Spot (\$) & DTE & ATM IV (\%) & Contracts \\
\midrule
NVDA & 207.04 & 80 & 44.6 & 142 \\
AMD  & 225.78 & 58 & 59.0 & 36 \\
MU   & 237.02 & 41 & 72.4 & 34 \\
INTC &  37.78 & 65 & 53.6 & 46 \\
\bottomrule
\end{tabular}
\end{table}

The calibrated parameters (Table~\ref{tab:supp_params}) produced the expected smile features: the negative $\beta_2$ generated put skew, the positive $\beta_4$ generated smile curvature, and the interaction term $\beta_3$ captured how the skew flattened at longer maturities. The calibration achieved an RMSE of 1.16\% IV on NVDA. To determine whether the smile shape generalized across tickers, we applied the $(\beta_1, \beta_2, \beta_3, \beta_4)$ calibrated on NVDA to AMD, MU, and INTC, re-optimizing only the scalar $\theta_{\text{base}}$ per ticker. Four parameters sufficed to capture the smile shape within the tradeable moneyness range, and the cross-ticker validation revealed that the shared smile shape transferred well to AMD (2.34\%) and reasonably to MU (3.99\%) and INTC (4.65\%; Table~\ref{tab:supp_rmse}). The higher RMSE for MU and INTC reflected differences in skew steepness across these names; MU exhibited a steeper put skew than NVDA, while INTC had a flatter OTM call wing, suggesting that per-ticker $\beta$ calibration would further reduce errors when multi-strike data was available. The residual analysis confirmed that the errors were concentrated in the wings ($K/S < 0.85$ and $K/S > 1.15$), with near-ATM residuals below 1\% IV for all tickers. These results established that the $\psi$-function captured the dominant smile features with a small number of parameters, and motivated the multi-expiration extension reported below.

\begin{table}[ht]
\centering
\caption{Calibrated smile parameters from the single-DTE NVDA fit (DTE = 80).}\label{tab:supp_params}
\begin{tabular}{llr}
\toprule
Parameter & Interpretation & Value \\
\midrule
$\beta_1$ & Term-structure decay (fixed) & 0.20 \\
$\beta_2$ & Skew (asymmetry)             & $-2.07$ \\
$\beta_3$ & DTE $\times$ skew interaction & 0.34 \\
$\beta_4$ & Smile curvature              & 1.43 \\
$\theta_{\text{base}}$ & Variance level (NVDA) & 0.084 \\
\bottomrule
\end{tabular}
\end{table}

\begin{table}[ht]
\centering
\caption{Cross-ticker RMSE (\% IV) from the single-DTE calibration. The smile shape $(\beta_2, \beta_3, \beta_4)$ was shared from NVDA; only $\theta_{\text{base}}$ was re-fitted per ticker.}\label{tab:supp_rmse}
\begin{tabular}{lrr}
\toprule
Ticker & $\theta_{\text{base}}$ (IV level) & RMSE (\% IV) \\
\midrule
NVDA (calibration) & 28.9\% & 1.16 \\
AMD  & 38.7\% & 2.34 \\
MU   & 49.0\% & 3.99 \\
INTC & 35.9\% & 4.65 \\
\bottomrule
\end{tabular}
\end{table}

\subsection{Multi-expiration NVDA calibration}

Having validated the smile shape across tickers, we addressed the term-structure limitation by collecting NVDA call and put chains at seven expirations on the same observation date (April 6, 2026): 0, 2, 4, 7, 14, 25, and 46 DTE, with NVDA trading near \$176.70. The cross-sectional snapshot provided the multi-expiration data needed to identify $\beta_1$ directly and to resolve the full term-structure geometry, yielding 73 observations after filtering to moneyness $K/S \in [0.90, 1.10]$ with positive volume. The six parameters $(\sigma_{\text{base}}, \beta_1, \ldots, \beta_5)$ of equation~\eqref{eq:psi} were fitted by minimizing mean squared IV error across all seven expirations simultaneously.

The calibrated coefficients (Table~\ref{tab:supp_term_params}) combined a negative $\beta_1$ with a positive $\beta_5$ to produce the U-shaped ATM curve: at short maturities the linear term dominated (IV decayed rapidly from the 0~DTE spike), while at longer maturities the quadratic term bent the curve upward. The trough occurred near DTE~$\approx 8$, where IV reached 32.5\%, a 1.38$\times$ ratio between the 0~DTE peak and the trough. The skew parameters ($\beta_2 = -8.45$, $\beta_3 = 1.62$) confirmed that the volatility skew steepened at shorter maturities and flattened as DTE increased, consistent with the concentration of crash risk premia in near-term options~\cite{bates2000}. The overall fit achieved an RMSE of 2.31\% IV across all seven expirations and an RMSE of 1.15\% IV across the 2 to 46 DTE range reported in the main text. The fit was tightest at intermediate and long maturities (7 to 46~DTE: RMSE $< 1.3$\%), with the largest residuals at 0~DTE (RMSE = 6.57\%) where the gamma-compressed smile shape deviated from the log-moneyness functional form. The 0~DTE level bias was eliminated ($-0.1$\%), confirming that the U-shaped term structure was captured correctly even though the extreme smile curvature at expiration remained a source of scatter.

\begin{table}[ht]
\centering
\caption{Calibrated $\psi$ parameters from the multi-expiration NVDA dataset.}\label{tab:supp_term_params}
\begin{tabular}{llr}
\toprule
Parameter & Interpretation & Value \\
\midrule
$\sigma_{\text{base}}$ & Base ATM IV level & 45.0\% \\
$\beta_1$ & DTE decay           & $-0.622$ \\
$\beta_2$ & Skew (asymmetry)    & $-8.448$ \\
$\beta_3$ & DTE $\times$ skew   & 1.617 \\
$\beta_4$ & Smile curvature     & $-1.242$ \\
$\beta_5$ & DTE curvature       & 0.150 \\
\bottomrule
\end{tabular}
\end{table}

\section{Sector-NN calibration: corpus partition and per-sector breakdown}\label{sec:supp_sector_nn}

The main text reports the headline three-way comparison (parametric vs.\ shared NN vs.\ sector NN) on the 31-ticker ladder. This section reports the corpus partition that defined the sector grouping and the per-sector and per-ticker decomposition that lay underneath the aggregate RMSE numbers. The corpus was dominated by the ETF and Technology sectors, which together carried two-thirds of the observation count (Table~\ref{tab:sector_map}); the per-sector decomposition revealed that the sector NN reduced RMSE in every sector over the shared NN, with the largest absolute gains in the sectors where the global shape was the worst fit (Table~\ref{tab:per_sector}); and the per-ticker fit ranking under the sector NN traced the residual error to a small set of names whose smile geometries pulled their sector RMSE upward (Table~\ref{tab:worst_tickers}).

\begin{table}[ht]
\centering
\caption{Sector partition of the 31-ticker ladder. Each ticker appeared on at least one of the two capture days (2026-04-14 and 2026-04-15); per-sector observation counts were pooled across days. ETF observation counts were substantially larger than other sectors because SPY, QQQ, and IWM carried far deeper strike ladders than individual-name chains.}\label{tab:sector_map}
\begin{tabular}{lrl}
\toprule
Sector & Observations & Tickers \\
\midrule
ETF        & 17{,}533 & IWM, QQQ, SPY \\
Technology &  7{,}689 & AAPL, AMD, AVGO, GOOG, INTC, META, MSFT, MU, NVDA, QCOM \\
Healthcare &  3{,}820 & ABBV, AMGN, BMY, JNJ, LLY, MRNA, PFE, UNH \\
Financials &  3{,}496 & BAC, GS, JPM, WFC \\
Retail     &  1{,}989 & TGT, UPS, WMT \\
Energy     &  1{,}434 & CVX, OXY, XOM \\
\bottomrule
\end{tabular}
\end{table}

\begin{table}[ht]
\centering
\caption{Per-sector RMSE (\% IV) across the three calibration strategies on the 31-ticker ladder. The sector NN reduced RMSE in every sector over the shared NN, with the largest absolute gains in sectors where the global shape poorly represented the local geometry.}\label{tab:per_sector}
\begin{tabular}{lrrrr}
\toprule
Sector & Parametric & Shared NN & Sector NN & Obs.\ count \\
\midrule
ETF        &  8.93 & 6.89 & 5.09 & 17{,}533 \\
Financials &  6.85 & 6.44 & 5.54 &  3{,}496 \\
Energy     &  7.66 & 6.44 & 5.97 &  1{,}434 \\
Retail     &  8.62 & 6.94 & 6.66 &  1{,}989 \\
Healthcare & 10.78 & 9.83 & 8.83 &  3{,}820 \\
Technology & 11.65 & 11.23 & 9.16 &  7{,}689 \\
\bottomrule
\end{tabular}
\end{table}

\begin{table}[ht]
\centering
\caption{Best and worst individual ticker fits under the sector-specific NN calibration, ranked by RMSE across all 31 tickers. Small-sample tickers (PFE, BMY) showed the expected small-sample penalty; the remaining worst fits were concentrated in technology and healthcare where within-sector shape heterogeneity was largest.}\label{tab:worst_tickers}
\begin{tabular}{llrrl}
\toprule
Rank & Ticker & Obs.\ count & RMSE (\% IV) & Sector \\
\midrule
\multicolumn{5}{l}{\emph{Best 5}} \\
 1 & GS   & 1{,}755 & 4.32 & Financials \\
 2 & WMT  &    676 & 4.71 & Retail \\
 3 & QQQ  & 6{,}386 & 5.06 & ETF \\
 3 & SPY  & 7{,}748 & 5.06 & ETF \\
 5 & XOM  &    498 & 5.09 & Energy \\
\midrule
\multicolumn{5}{l}{\emph{Worst 5}} \\
27 & JNJ  &    408 & 11.79 & Healthcare \\
28 & ABBV &    433 & 12.18 & Healthcare \\
29 & INTC &    531 & 12.47 & Technology \\
30 & AVGO &    594 & 12.59 & Technology \\
31 & PFE  &    101 & 13.33 & Healthcare \\
\bottomrule
\end{tabular}
\end{table}

\section{Holdout diagnostics motivating the earnings-aware mechanism}\label{sec:supp_earnings_diag}

The temporal-holdout result of the main text attributed the test-time generalization gap to scheduled corporate events rather than to a broad-market regime shift, but that attribution required ruling out the alternative explanation that the 04-23 and 04-24 capture days landed inside a volatility expansion that the training block did not see. Two diagnostics established the alternative as untenable and pinned the blow-up to ticker-level event behavior. First, the broad-market volatility index VIX traced a flat path through the holdout window: VIX closed at comparable levels on the final training day and the two holdout days, with no intraday excursion during the holdout that approached even a one-standard-deviation move relative to the trailing six-day window. A regime-shift hypothesis required the index to move materially over the holdout, and it did not. Second, the per-ticker ATM~IV trace revealed that the IV expansion on 04-23 and 04-24 was concentrated entirely in the technology sector, and within that sector concentrated on the names with earnings prints inside or adjacent to the holdout window: INTC (which printed inside the window), MSFT, META, and GOOG all exhibited ATM~IV elevations of 5 to 30 percentage points relative to their training-block baselines, while non-Tech names traced ATM~IV paths consistent with their pre-holdout levels. The combination established that the holdout RMSE jump was not a market-wide effect that the architecture had failed to absorb but a localized response to the earnings calendar, and it directly motivated the earnings indicator and peer-coupling feature reported in Configuration~C of Table~\ref{tab:earnings_holdout}. The per-ticker decomposition of the earnings-aware fit on the Tech sector (Table~\ref{tab:tech_pivot}) traced the Configuration~C improvement to the peer-coupling feature: tickers whose own prints fell outside the holdout but whose sector peers printed during it received the largest test-RMSE reductions, while INTC, the print-day ticker, was the one Tech name to worsen because its $+2192\%$ EPS surprise produced a post-print IV crush that no calendar-based indicator could anticipate.

\begin{table}[ht]
\centering
\caption{Per-ticker test RMSE on the Tech sector across the pooled control (A) and the earnings-aware fit (C). Columns $e$ and $e^{\text{peer}}$ report the median signed days-to-earnings and the minimum same-sector peer distance for each ticker on the holdout days. The largest gains accrued to tickers whose peers printed within the holdout (QCOM, AMD, META, MSFT, GOOG); INTC, the print-day ticker, was the one name that worsened, reflecting an unanticipated outcome that no calendar-based indicator could absorb.}\label{tab:tech_pivot}
\begin{tabular}{lrrrrr}
\toprule
Ticker & $e$ & $e^{\text{peer}}$ & A: pooled (\%) & C: earnings-aware (\%) & $\Delta$ (\%) \\
\midrule
QCOM & $+5$  & $1$ & 41.66 & 37.59 & $-4.07$ \\
AMD  & $+11$ & $1$ & 40.07 & 36.50 & $-3.57$ \\
META & $+5$  & $1$ & 14.32 & 10.73 & $-3.59$ \\
MSFT & $+5$  & $1$ & 14.00 & 10.73 & $-3.28$ \\
GOOG & $+5$  & $1$ &  9.87 &  7.90 & $-1.97$ \\
MU   & $-30$ & $1$ & 16.95 & 15.59 & $-1.36$ \\
AAPL & $+6$  & $1$ & 10.42 & 10.34 & $-0.08$ \\
AVGO & $+30$ & $1$ & 12.58 & 12.83 & $+0.24$ \\
NVDA & $+26$ & $1$ &  9.43 & 10.47 & $+1.05$ \\
INTC & $-1$  & $5$ & 33.17 & 35.94 & $+2.77$ \\
\bottomrule
\end{tabular}
\end{table}

\section{LLY short-premium scenario: terminal P\&L statistics}\label{sec:supp_lly_pnl}

The main-text §5 short-premium scenario rendered the share-price, premium, IV, Greek, and terminal-P\&L distributions for the LLY 31-day forward simulation as figures. The corresponding moments and tail statistics per contract are reported in Table~\ref{tab:lly_scenario}. Both legs showed the short-premium asymmetry that the figures rendered qualitatively: a thin-tailed right side with a hard ceiling at the entry premium, paired with a long left tail that extended an order of magnitude beyond the median. The model's $t=0$ Leisen-Reimer fair value sat above the market mid on both contracts (entry edge $+\$2.61$ on the put, $+\$0.58$ on the call), so the shorts were sold at a small discount to the model. The dispersion was substantially larger on the call leg (std $\$42.36$ vs.\ $\$20.22$), and the worst-case loss on the call ($-\$325.62$) exceeded the worst put loss ($-\$174.60$) by a factor of $1.87$, despite the call collecting a smaller entry premium.

\begin{table}[ht]
\centering
\caption{Terminal P\&L statistics from 1{,}000 LLY 31-day forward simulations of the real 2026-05-29 \$825 put and \$940 call. Entry premium was the market mid from the 2026-04-28 capture. The model's $t=0$ Leisen-Reimer fair value at NN-IV was reported as a reference; statistics are per contract.}
\label{tab:lly_scenario}
\begin{tabular}{lrr}
\toprule
Statistic & Short put (\$) & Short call (\$) \\
\midrule
Strike $K$                  & 825        & 940       \\
Market mid (entry premium)  & $+23.30$   & $+20.76$  \\
Model $t=0$ fair value      & $+25.91$   & $+21.34$  \\
Entry edge (model $-$ market) & $+2.61$  & $+0.58$   \\
\midrule
Mean P\&L                    & $+16.94$  & $+1.86$   \\
Median P\&L                  & $+23.30$  & $+20.76$  \\
Std P\&L                     &  $20.22$  & $42.36$   \\
5\%-tile P\&L                & $-28.47$  & $-97.56$  \\
Worst-case P\&L              & $-174.60$ & $-325.62$ \\
Premium kept in full         & $83.1\%$  & $70.2\%$  \\
\bottomrule
\end{tabular}
\end{table}
